\textbf{\textit{Soal.}} Diberikan segitiga $ABC$ dan lingkaran luar $\omega$ dari segitiga tersebut. Ditarik sembarang garis yang sejajar $BC$, yang memotong sisi $AB$ dan $AC$ secara berurutan di titik $X$ dan $Y$. Misalkan $M$ adalah titik tengah dari busur $BC$ pada lingkaran $\omega$ yang tidak memuat titik $A$. Garis $MX$ dan $MY$ memotong lingkaran $\omega$ kembali di titik $S$ dan $T$ secara berurutan. Misalkan $ST$ dan $XY$ berpotongan di titik $P$. Tunjukkan bahwa garis $AP$ menyinggung lingkaran $\omega$.


\begin{proof}[\textbf{Solusi. }] (Senin, 7 April 2025) 
    \begin{center}
        \begin{asy}
         /* Geogebra to Asymptote conversion, documentation at artofproblemsolving.com/Wiki go to User:Azjps/geogebra */
        import graph; size(7.475886785948781cm); 
        real labelscalefactor = 1; /* changes label-to-point distance */
        pen dps = linewidth(0.7) + fontsize(10); defaultpen(dps); /* default pen style */ 
        pen dotstyle = black; /* point style */ 
        real xmin = -3.855436869028034, xmax = 3.6204499169207467, ymin = -3.793232734663582, ymax = 1.5347953828441994;  /* image dimensions */
        pen rvwvcq = rgb(0.08235294117647059,0.396078431372549,0.7529411764705882); pen wewdxt = rgb(0.43137254901960786,0.42745098039215684,0.45098039215686275); 
        
        draw((-0.09740583514552568,0.26281957034906533)--(-0.6502960977693696,-2.544476144622696)--(3.029372008454863,-2.4945258810210604)--cycle, linewidth(1.) + rvwvcq); 
         /* draw figures */
        draw((-0.09740583514552568,0.26281957034906533)--(-0.6502960977693696,-2.544476144622696), linewidth(1.) + rvwvcq); 
        draw((-0.6502960977693696,-2.544476144622696)--(3.029372008454863,-2.4945258810210604), linewidth(1.) + rvwvcq); 
        draw((3.029372008454863,-2.4945258810210604)--(-0.09740583514552568,0.26281957034906533), linewidth(1.) + rvwvcq); 
        draw(circle((1.1749636876576037,-1.4458636762846033), 2.130380867329792), linewidth(1.) + wewdxt); 
        draw((-0.09740583514552568,0.26281957034906533)--(-2.4196935987651527,-1.466470032875881), linewidth(1.) + wewdxt); 
        draw((1.7986414439874552,-1.4092075458083395)--(-2.4196935987651527,-1.466470032875881), linewidth(1.) + wewdxt); 
        draw((-2.4196935987651527,-1.466470032875881)--(2.287296805328559,0.37106716313038374), linewidth(1.) + wewdxt); 
        draw((-0.8741658808555909,-0.8631196061047632)--(1.2038802311409873,-3.576048286452888), linewidth(1.) + wewdxt); 
        draw((2.287296805328559,0.37106716313038374)--(1.2038802311409873,-3.576048286452888), linewidth(1.) + wewdxt); 
        draw((-0.09740583514552568,0.26281957034906533)--(1.2038802311409873,-3.576048286452888), linewidth(1.) + wewdxt); 
         /* dots and labels */
        dot((-0.09740583514552568,0.26281957034906533),dotstyle); 
        label("$A$", (-0.06754323026779428,0.3443141003385546), NW * labelscalefactor); 
        dot((-0.6502960977693696,-2.544476144622696),dotstyle); 
        label("$B$", (-0.6169959339121367,-2.461225705286637), NE * labelscalefactor); 
        dot((3.029372008454863,-2.4945258810210604),dotstyle); 
        label("$C$", (3.062672172312096,-2.4112754416850013), SE * labelscalefactor); 
        dot((-0.4326730432884509,-1.4394968935430041),dotstyle); 
        label("$X$", (-0.400544868840123,-1.3539948621170508), NE * labelscalefactor); 
        dot((1.7986414439874552,-1.4092075458083395),linewidth(4.pt) + dotstyle); 
        label("$Y$", (1.8305661095944796,-1.3456698181834448), E * labelscalefactor); 
        dot((1.2038802311409873,-3.576048286452888),linewidth(4.pt) + dotstyle); 
        label("$M$", (1.2394882011285961,-3.510181240920981), NE * labelscalefactor); 
        dot((-0.8741658808555909,-0.8631196061047632),linewidth(4.pt) + dotstyle); 
        label("$S$", (-0.8417720399484586,-0.7962169185654548), NW * labelscalefactor); 
        dot((2.287296805328559,0.37106716313038374),linewidth(4.pt) + dotstyle); 
        label("$T$", (2.3217435264886643,0.4358895836082196), NE * labelscalefactor); 
        dot((-2.4196935987651527,-1.466470032875881),linewidth(4.pt) + dotstyle); 
        label("$P$", (-2.3902296593097874,-1.4039451257186861), NW * labelscalefactor); 
        dot((0.4027680784895119,-2.5301811507060994),linewidth(4.pt) + dotstyle); 
        label("$D$", (0.4319592275906988,-2.461225705286637), NE * labelscalefactor); 
        clip((xmin,ymin)--(xmin,ymax)--(xmax,ymax)--(xmax,ymin)--cycle); 
         /* end of picture */
    \end{asy}
\end{center}
    Perhatikan bahwa $\triangle AXY \sim \triangle ABC$ sehingga $\angle AYX = \angle ACB$. Lalu karena $M$ titik tengah busur $BC$, maka $BM=MC$ dan $\angle BCM = \angle MBC$.
    \begin{lemma}
        $STYX$ merupakan segiempat siklis.
    \end{lemma}
    \begin{proof}[\textbf{Bukti 1.}]
        Misalkan $D$ adalah perpotongan $XM$ dan $BC$.
       \begin{align*}
        \dangle STY &= \dangle STM\\
        &= \dangle STB + \dangle BTM\\
        &= \dangle SMB + \dangle BCM\\
        &= \dangle XMB + \dangle MBC\\
        &= \dangle MDB\\
        &= \dangle XDB\\
        &= \dangle XBC + \dangle MXB\\
        &= \dangle AXY + \dangle SXA\\
        &= \dangle SXY.
       \end{align*}
       Terbukti.
    \end{proof} 
    \begin{proof}[\textbf{Bukti 2.}]
        Kita punya $\dangle XSB = \dangle MSB = \dangle MCB = \dangle CBM$. Dari sini diperoleh
        \begin{align*}
            \dangle SXY &= \dangle SXA + \dangle AXY\\
            &= \dangle MXB + \dangle ABC\\
            &= \dangle SXB + \dangle ABC\\
            &= \dangle SBX + \dangle XSB + \dangle ABC\\
            &= \dangle SBX + \dangle CBM + \dangle ABC\\
            &= \dangle SBM\\
            &= \dangle STM\\
            &= \dangle STY.
        \end{align*}
        Terbukti.
    \end{proof}
    \begin{proof}[\textbf{Bukti 3.}]
        Akan dibuktikan bahwa $MY \cdot MT = MX \cdot MS$. Perhatikan dari P.O.P bahwa
        \begin{align*}
            AY \cdot YC &= MY \cdot YT\\
            AX \cdot XB &= MX \cdot XS.
        \end{align*}
        Berarti karena $AY=YC$ dan $AX=XB$, maka 
        \begin{align*}
            MY \cdot MT &= MY(MY+YT) = MY^2 + MY \cdot YT = MY^2 + AY^2\\
            MX \cdot MS &= MX(MX+XS) = MX^2 + MX \cdot XS = MX^2 + AX^2.
        \end{align*}
        Lalu dengan dalil cosinus pada $\triangle AMB$ dan $\triangle AMC$, dan $\cos B = \cos(180- C) = -\cos C$, kita peroleh
        \begin{align*}
            AM^2 &= MC^2 + AC^2 - 2MC \cdot AC \cdot \cos C\\
            AM^2 &= MB^2 + AB^2 - 2MB \cdot BA \cdot \cos B = MC^2 + AB^2 + 2MC \cdot BA \cdot \cos C.
        \end{align*}
        sehingga dengan mengeliminasi kedua persamaan tersebut didapat
        \begin{align*}
            AC^2 - AB^2 &= 2MC (AC+AB) \cos C\\
            AC - AB &= 2MC \cos C.
        \end{align*}
        Lalu dengan dalil cosinus pada $\triangle BXM$ dan $\triangle CYM$, kita peroleh
        \begin{align*}
            MY^2 &= AY^2 + MC^2 - 2MC \cdot AY \cdot \cos C\\
            MX^2 &= AX^2 + MB^2 - 2MB \cdot AX \cdot \cos B = AX^2 + MC^2 - 2MC \cdot AX \cdot \cos C.
        \end{align*}
        sehingga dengan mengeliminasi kedua persamaan tersebut didapat
        \begin{align*}
            &MY \cdot MT - MX \cdot MS\\ 
            &= (MY^2 + AY^2) - (MX^2 + AX^2)\\
            &= (MY^2 - MX^2) + (AY^2 - AX^2)\\
            &= (AY^2 + MC^2 - 2MC \cdot AY \cdot \cos C) - (AX^2 + MC^2 - 2MC \cdot AX \cdot \cos C) + (AY^2 - AX^2)\\
            &= 2(AY^2 - AX^2) - 2 MC \cdot (AX+AY) \cdot \cos C\\
            &= \frac12((2AY)^2 - (2AX)^2) + \frac12 \cdot 2 MC \cos C \cdot 2(AX+AY) \\
            &= \frac12(AB^2-AC^2) + \frac12(AC-AB)(AB+AC)\\
            &= \frac12(AB^2-AC^2) + \frac12(AC^2-AB^2)\\
            &= 0.
        \end{align*}
        yang menyebabkan $MY \cdot MT = MX \cdot MS$. Terbukti.
    \end{proof} 
    Sekarang perhatikan bahwa \textit{radical axis} dari lingkaran $(STYX)$ dan $\omega$ adalah $ST$. 
    Lalu \textit{radical axis} dari lingkaran $(STYX)$ dan $(AXY)$ adalah $XY$. 
    Terakhir, \textit{radical axis} dari $\omega$ dan $(AXY)$ adalah garis singgung yang melewati $A$ (yang menyinggung dua lingkaran tersebut dari \textit{alternate segment theorem}).
    Oleh karena itu, $ST$, $XY$, dan garis singgung yang melewati $A$ berpotongan di satu titik \textit{radical center}, yaitu $P$. Terbukti.
\end{proof} 